\chapter{理论研究}

此处格式已按模板设定，作者只需选择段落区域，输入替换之。模板中所有说明性文字用于注释格式与内容的要求，撰写论文时请删除。模板中，图表、公式、参考文献等都已给出范例，撰写论文时请删除。

本模板已包含符合章节设置的“多级别列表”，只需在相应位置替换标题文字即可。如需增加章节，建议先使用格式刷，再调整编号。

\section{毕业设计（论文）结构及要求}

\subsection{论文结构}

毕业设计（论文）原则上应采用中文完成（教学语言为外语的除外），确需用其他语言撰写的需经学院审批同意并报教务处备案。毕业设计（论文）一般由以下部分组成，依次为：\ding{172}封面；\ding{173}扉页；\ding{174}独创性声明；\ding{175}中英文摘要及关键词；\ding{176}目录；\ding{177}正文；\ding{178}参考文献；\ding{179}附录；\ding{180}致谢。

\subsection{语言表述}

要做到数据可靠、推理严谨、立论正确。论述必须简明扼要、重点突出，对同行专业人员已熟知的常识性内容，尽量减少叙述。

论文中如出现一些非通用性的新名词、术语或概念，需做出解释。

\subsection{标题和层次}

标题要重点突出，简明扼要，层次清晰。

\subsection{打印规格}

论文一律采用A4纸张（大小为210mm$\times$297mm）单面打印，页边距如下设置。上：27.5mm；下：25.4mm；左：35.7mm；右：27.7mm。页眉距边界15.0mm；页脚距边界17.5mm。字符间距为默认值（缩放100\%，间距：标准）。

\section{毕业设计（论文）撰写规范}

\subsection{封面}

采用天津大学本科生毕业设计（论文）统一封面，封面内容包括论文题目、学院、专业、年级、姓名、学号、指导教师等信息，居中编排。

论文题目是论文总体内容的体现，要求醒目、简明、准确，主题突出，一般不宜超过25字。

\subsection{目录}

目录的各章节应简明扼要，应列至三级标题，包含正文及其后的各部分，并附有相应页码。目录的文字应与相应标题文字完全一致。

“目录”两字之间空一个全角空格或两个半角空格，采用不编号章标题样式。目录条目采用正文样式。
各级标题采用逐级缩进形式，每级缩进2字符，页码前导符采用“…”。

\subsection{正文}

正文是毕业设计（论文）的主体，应占据主要篇幅，文字一般不少于15000字，要求主题明确，内容充实；论点正确，论据可靠，论证充分。内容一般包括：设计（论文）的工作目的（背景），国内外研究现状、理论分析、计算方法（研究方法）、实验装置和测试方法、实验结果分析与讨论、研究成果、结论及意义等。

正文中文字体为宋体，英文字体为Times New Roman，正文采用小四号字，段落行间距为固定值20磅，段落前后间距为0，首行缩进2字符。西文字体以Times New Roman为准，若Times New Roman中没有相应字符，则应使用较为清晰和通用的字体，语料文字的字体可采用楷体。文本语料、数学公式和专门文字（如计算机程序代码）的字体可以根据需要选择。

章节与标号：一般分为章标题（一级标题）、不编号章标题（同属于一级标题）、二级标题和三级标题。各章节编号建议采用Word的“多级别列表”方式自动形成编号，标题编号与标题内容之间空一个全角空格或两个半角空格。各级章节标题格式要求细节参见《天津大学本科生毕业设计（论文）相关材料撰写规范》。

图、表等与其前后的正文之间要有一行的间距；文中的图、表、公式一律采用阿拉伯数字分章编号，如：图2-5，表3-2，公式（5-1）（“公式”二字不出现）等。若图或表中有附注，采用英文小写字母顺序编号。子图采用英文字母编号。引用图或表应在图题或表题右上角标出文献来源。图或表的附注应位于图或表的下方。

\subsection{公式}

公式要标准、通用，公式的变量和参数，除了公知公认的以外，一般要给出解释。文中的公式一律采用阿拉伯数字分章编号，如：公式（5-1）（“公式”二字不出现）。公式编号只能用于行间公式。公式编号右对齐，且应加英文小括号，不加引导符。其余样式与正文相同。变量和参数，不能采用正文格式，必须正确使用数学格式或行内公式格式，包含行内公式的段落可以采用单倍行距。

依照以上标准的行间公式范例如下。
\begin{equation}
    \begin{cases}
        F_{si} &= k_{si}(z'_i-z_i)+c_{si}(\dot{z}'_i-\dot{z}_i) \\
        F_{ti} &= k_{ti}(z_i-y_i(x_i,t))+c_{ti}(\dot{z}_i-\dot{y}_i(x_i,t))
    \end{cases}
    \label{equ:\theequation}
\end{equation}

其中：$z_i$为车辆第i个轮胎由静平衡位置起算的竖向位移；$y_i$为桥梁在第i个轮胎作用下的瞬时变位。

\subsection{图}

图要精选、简明，切忌与表及文字表述重复。图中的术语、符号、单位等应同文字表述一致。图与其前后的正文之间要有一行的间距；文中的图一律采用阿拉伯数字分章编号，如：图2-5。若图中有附注，采用英文小写字母顺序编号。子图采用英文字母编号。图序及图题居中置于图的下方，图、表内容为五号字，中文字体为宋体，英文字体为Times New Roman。图序与图题之间空一个全角空格或两个半角空格。引用图应在图题右上角标出文献来源。图的附注应位于图或表的下方。

依照以上标准的插图范例如下。

\figuremacro{htbp!}
            {晶体坐标系和样品坐标系示意图.jpg}
            {晶体坐标系和样品坐标系示意图}
            {0.5}
            {晶体}

\figuremacro{htbp!}
            {应变硅横截面样品的拉曼类硅峰峰位.png}
            {应变硅横截面样品的拉曼类硅峰峰位 \\ (a)  云图；(b)  沿深度方向分布曲线}
            {0.8}
            {应变}

\subsection{表}

表中参数应标明量和单位的符号。表的编排建议采用国际通行的三线表。表与其前后的正文之间要有一行的间距；文中的表一律采用阿拉伯数字分章编号，如：表3-2。若表中有附注，采用英文小写字母顺序编号。表序及表题居中置于表的上方。如某个表需要转页接排，在随后的各页上应重复表序，后跟表题（可省略）和“（续）”，置于表格上方。续表均应重复表头。表序与表题之间空一个全角空格或两个半角空格。引用表应在表题右上角标出文献来源。表的附注应位于图或表的下方。

依照以上标准的三线表范例如下。

\tablemacro{htbp}
            {典型微尺度力学实验方法的基本信息}
            {ccc}
            {\songtibf{实验方法} & \songtibf{主要测量对象} & \songtibf{空间分辨率} \\}
            {原子力显微镜 & 表面力 & $0.1 nm$\cite{Zhang2002} \\
            透射电镜 & 晶格结构、位错 & $0.1 nm$\cite{Borko1978} \\
            X射线衍射 & 应变 & $1 \mu m$\cite{Hao2008,Dupont1974,Wang1997} \\
            同步辐射 & 内部三维结构与变形 & 约$100 nm$\cite{Liang2024,Koya1995} \\
            显微拉曼 & 应变 & $250 nm$\cite{Deng2006} \\}
            {基本信息}

\clearpage %为了与原模板分页保持一致添加，实际使用不需要

\section{脚注}

在正文当中\footnote{脚注字号为五号，其余样式与正文相同。}，如果有个别名词或情况需要解释时，可加注释说明。注释说明应采用文中编号加脚注的模式。脚注应采用阿拉伯数字上标，字体为Times New Roman，分页连续标号。脚注字号为五号，其余样式与正文相同。

\section{页眉和页脚}

页眉从正文开始到毕业设计（论文）结尾，一律设为“天津大学20XX届本科生毕业设计（论文）”，采用五号字居中编排，其余样式与正文相同。

页脚从中文摘要开始，从Ⅰ开始顺序编页码，为大写罗马数字格式。正文第一页开始，用阿拉伯数字重新从1开始顺序编页码，至学位论文结尾。页脚的字号为小五号，居中编排。

\clearpage